%% =========================================================================
\section{Introduction}
%% =========================================================================

This report presents a rigorous simulation-based investigation of three-phase
synchronous generator behaviour using MATLAB R2025b and Simscape Electrical.
The study is structured across three exercises and six tasks of increasing
complexity.

Exercise~1 analyses the dynamic response of a \MVA{555}, \kV{24}, \Hz{60}
round-rotor generator to step load changes, establishing the governing
relationships between active power and system frequency (described by the
swing equation) and between reactive power and terminal voltage (governed by
the Automatic Voltage Regulator, AVR). Exercise~2 extends the study to the
parallel operation of two identical synchronous generators, demonstrating that
the governor frequency setpoint determines active power sharing: an
under-frequency machine absorbs power and operates as a synchronous motor,
while an over-frequency machine supplies the load. The system frequency settles
at a droop-weighted average of both setpoints. Exercise~3 replaces one
generator with an ideal infinite bus, revealing the fundamental distinction from
the two-machine case: the infinite bus holds frequency absolutely rigid at
\pu{1.0}, so the oncoming generator cannot influence system frequency, and its
power output is determined entirely by its governor setpoint relative to the
fixed grid frequency. All results are supported by simulation plots, governing
equation derivations, and quantitative steady-state analysis.
